% ==============================================
% 第6章 総仕上げ模試
% ==============================================
\chapter{総仕上げ模試}

いよいよ最終章です。ここまで学んできたストラテジ系・マネジメント系・テクノロジ系の知識を、模試形式で総点検しましょう。

\begin{keypoint}
この章の構成：
\begin{itemize}
  \item \textbf{6.1 分野別ミニ模試} — 分野ごとに5問ずつ。苦手分野の発見に
  \item \textbf{6.2 本試験形式模試} — 3分野混合で本番と同じ流れを体験
  \item \textbf{6.3 公式問題セクション} — 過去問を追加するための受け皿
  \item \textbf{6.4 最終チェックリスト} — 試験直前の確認用
\end{itemize}
\end{keypoint}

\begin{stumble}
模試は「点数」ではなく「間違えたテーマ」を確認するためのもの。間違えた問題の解説を読み、該当する章に戻って復習するのが最も効率的な使い方です。
\end{stumble}

% ====================================================
%  6.1 分野別ミニ模試
% ====================================================
\section{分野別ミニ模試}

まずは分野ごとに5問ずつ解いて、弱点を洗い出しましょう。各ミニ模試の目安時間は\textbf{8分}です。

% -----------------------------------------------
%  ストラテジ系ミニ模試
% -----------------------------------------------
\mondaipage
\begin{minimoshi}{ストラテジ系}

制限時間 \textbf{8分} で以下の5問を解いてください。

\begin{mondai}{企業が自社の経営環境を分析するために、内部環境の強み・弱みと外部環境の機会・脅威を整理するフレームワークはどれか。}{basic}
  \choice{BSC（バランススコアカード）}
  \choice{PPM（プロダクト・ポートフォリオ・マネジメント）}
  \choice{SWOT分析}
  \choice{バリューチェーン分析}
\end{mondai}

\begin{mondai}{DX（デジタルトランスフォーメーション）の説明として最も適切なものはどれか。}{standard}
  \choice{紙の書類を電子化すること}
  \choice{社内のパソコンを最新機種に入れ替えること}
  \choice{デジタル技術を活用してビジネスモデルや組織を変革すること}
  \choice{Webサイトを開設してオンライン販売を始めること}
\end{mondai}

\begin{mondai}{損益分岐点の説明として最も適切なものはどれか。}{standard}
  \choice{売上高が最大になる点}
  \choice{売上高と総費用が等しくなり、利益がゼロになる点}
  \choice{固定費がゼロになる点}
  \choice{変動費率が100\%になる点}
\end{mondai}

\begin{mondai}{個人情報保護法において、個人情報に該当しないものはどれか。}{trick}
  \choice{氏名}
  \choice{生年月日と氏名の組合せ}
  \choice{会社の代表電話番号}
  \choice{顔写真}
\end{mondai}

\begin{mondai}{SaaS の説明として最も適切なものはどれか。}{standard}
  \choice{サーバやネットワークなどのインフラをインターネット経由で提供するサービス}
  \choice{アプリケーション開発のためのプラットフォームをインターネット経由で提供するサービス}
  \choice{アプリケーションの機能をインターネット経由で提供するサービス}
  \choice{データベースの管理機能だけを提供するサービス}
\end{mondai}

\end{minimoshi}

\kaitoupage
\subsection*{ストラテジ系ミニ模試 解答}

\begin{enumerate}
  \item \textbf{C}（SWOT分析）— 内部環境（強み・弱み）と外部環境（機会・脅威）の2軸4象限で整理するのがSWOT分析。BSCは4つの視点で経営を評価、PPMは市場成長率と占有率で事業を分類、バリューチェーンは付加価値の源泉を分析する。
  \item \textbf{C}（ビジネスモデルや組織の変革）— DXは単なるIT化ではなく、デジタル技術によって\textbf{ビジネスモデルそのもの}を変革すること。A は単純なデジタイゼーション、B は設備更新、D はEC開設にすぎない。
  \item \textbf{B}（売上高＝総費用の点）— 損益分岐点は「利益がゼロ」になるポイント。これを超えると黒字、下回ると赤字。損益分岐点売上高＝固定費÷（1－変動費率）で計算できる。
  \item \textbf{C}（会社の代表電話番号）— 会社の代表電話番号は特定の個人を識別できないため個人情報に該当しない。氏名・生年月日と氏名の組合せ・顔写真はいずれも個人を識別できる情報。
  \item \textbf{C}（アプリケーション機能の提供）— SaaS（Software as a Service）はアプリケーションそのものを提供。A はIaaS、B はPaaS。「SaaS → ソフト、PaaS → 開発基盤、IaaS → インフラ」と覚える。
\end{enumerate}

% -----------------------------------------------
%  マネジメント系ミニ模試
% -----------------------------------------------
\mondaipage
\begin{minimoshi}{マネジメント系}

制限時間 \textbf{8分} で以下の5問を解いてください。

\begin{mondai}{WBS（Work Breakdown Structure）の説明として最も適切なものはどれか。}{basic}
  \choice{プロジェクトの作業を階層的に分解した構成図}
  \choice{プロジェクトのリスクを一覧にした表}
  \choice{プロジェクトの進捗を棒グラフで表した図}
  \choice{プロジェクトのコストを時系列で表したグラフ}
\end{mondai}

\begin{mondai}{アジャイル開発の特徴として最も適切なものはどれか。}{standard}
  \choice{要件定義から運用まで順番に工程を進め、前の工程には戻らない}
  \choice{短いイテレーション（反復）を繰り返して段階的に機能を追加する}
  \choice{プロトタイプを作成して顧客に確認し、一度で完成させる}
  \choice{外部の開発会社にすべてを委託し、完成品を受け取る}
\end{mondai}

\begin{mondai}{ITIL における「インシデント管理」の目的として最も適切なものはどれか。}{standard}
  \choice{障害の根本原因を究明し、再発を防止すること}
  \choice{ITサービスの中断を可能な限り迅速に復旧すること}
  \choice{ITサービスの変更を安全に実施すること}
  \choice{ITサービスの構成情報を正確に管理すること}
\end{mondai}

\begin{mondai}{SLA（Service Level Agreement）の説明として最も適切なものはどれか。}{standard}
  \choice{システム開発の作業手順を定めた文書}
  \choice{サービス提供者と利用者の間でサービスの品質水準を合意した文書}
  \choice{情報セキュリティに関するポリシーを定めた文書}
  \choice{プロジェクトの予算と期間を定めた契約書}
\end{mondai}

\begin{mondai}{システム監査人に求められる要件として最も適切なものはどれか。}{trick}
  \choice{監査対象のシステムの開発に参加した経験があること}
  \choice{監査対象から独立した立場であること}
  \choice{監査対象部門の上長の指示に従うこと}
  \choice{監査結果を監査対象部門だけに報告すること}
\end{mondai}

\end{minimoshi}

\kaitoupage
\subsection*{マネジメント系ミニ模試 解答}

\begin{enumerate}
  \item \textbf{A}（作業を階層的に分解した構成図）— WBSはプロジェクトの作業を「大→中→小」と分解し、全体を見渡せるようにする。B はリスク登録簿、C はガントチャート、D はEVMなど。
  \item \textbf{B}（短い反復で段階的に開発）— アジャイル開発は1〜4週間のイテレーションを繰り返す。A はウォーターフォール開発、C はプロトタイピング、D はアウトソーシング。
  \item \textbf{B}（迅速な復旧）— インシデント管理は「まず復旧」が目的。A は問題管理（根本原因の究明）、C は変更管理、D は構成管理。「インシデント＝応急処置、問題＝根治」と覚える。
  \item \textbf{B}（サービスの品質水準の合意文書）— SLAはサービスレベル（稼働率99.9\%など）を取り決める。A は作業手順書、C はセキュリティポリシー、D は契約書の一般的な内容。
  \item \textbf{B}（監査対象から独立した立場）— システム監査の最重要原則は\textbf{独立性}。A は独立性に反する（自分が作ったシステムを自分で監査してはならない）。C は独立性を損なう。D は経営者に報告するのが正しい。
\end{enumerate}

% -----------------------------------------------
%  テクノロジ系ミニ模試
% -----------------------------------------------
\mondaipage
\begin{minimoshi}{テクノロジ系}

制限時間 \textbf{8分} で以下の5問を解いてください。

\begin{mondai}{情報セキュリティの3要素（CIA）に含まれるものの組合せとして正しいものはどれか。}{basic}
  \choice{機密性・完全性・責任追跡性}
  \choice{機密性・完全性・可用性}
  \choice{機密性・信頼性・可用性}
  \choice{否認防止・完全性・可用性}
\end{mondai}

\begin{mondai}{公開鍵暗号方式の説明として最も適切なものはどれか。}{standard}
  \choice{暗号化と復号に同じ鍵を使用する方式}
  \choice{暗号化に公開鍵、復号に秘密鍵を使用する方式}
  \choice{鍵を使わずにデータを暗号化する方式}
  \choice{通信のたびに鍵を自動的に変更する方式}
\end{mondai}

\begin{mondai}{関係データベースにおいて、表の行を一意に識別するために設定するものはどれか。}{basic}
  \choice{外部キー}
  \choice{インデックス}
  \choice{主キー}
  \choice{ビュー}
\end{mondai}

\begin{mondai}{TCP/IP のネットワーク層のプロトコルはどれか。}{standard}
  \choice{HTTP}
  \choice{TCP}
  \choice{IP}
  \choice{SMTP}
\end{mondai}

\begin{mondai}{フィッシング攻撃の説明として最も適切なものはどれか。}{trick}
  \choice{大量のデータを送信してサーバをダウンさせる攻撃}
  \choice{実在する組織を装ったメールやWebサイトで個人情報を騙し取る攻撃}
  \choice{ネットワークを流れるデータを盗聴する攻撃}
  \choice{ソフトウェアの脆弱性を悪用して不正プログラムを実行する攻撃}
\end{mondai}

\end{minimoshi}

\kaitoupage
\subsection*{テクノロジ系ミニ模試 解答}

\begin{enumerate}
  \item \textbf{B}（機密性・完全性・可用性）— CIAはConfidentiality（機密性）・Integrity（完全性）・Availability（可用性）。責任追跡性・信頼性・否認防止は7要素に含まれるが、基本の3要素ではない。
  \item \textbf{B}（公開鍵で暗号化、秘密鍵で復号）— 公開鍵暗号方式は暗号化と復号で異なる鍵を使う。A は共通鍵暗号方式の説明。公開鍵は「南京錠」、秘密鍵は「その鍵」にたとえると分かりやすい。
  \item \textbf{C}（主キー）— 主キーは表の各行を一意に識別するフィールド。A の外部キーは他の表との関連付け、B のインデックスは検索高速化、D のビューは仮想表。
  \item \textbf{C}（IP）— TCP/IPモデルのネットワーク層（インターネット層）で動作するのはIP。A のHTTPとD のSMTPはアプリケーション層、B のTCPはトランスポート層。
  \item \textbf{B}（偽サイト等で個人情報を騙し取る）— フィッシングは「釣り」の意味。A はDoS攻撃、C は盗聴（スニッフィング）、D はエクスプロイト（脆弱性悪用攻撃）。
\end{enumerate}

% ====================================================
%  6.2 本試験形式模試
% ====================================================
\section{本試験形式模試}

本試験は \textbf{100問・120分} です。合格基準は\textbf{総合600点以上}かつ\textbf{各分野300点以上}（1000点満点、IRT方式）。

ここでは本番と同じ「分野混合」形式で出題します。分野をまたいで考える力を確認してください。

\vspace{0.5em}
\noindent ※ 実際の試験は100問です。本模試では代表的な20問を収録しています。過去問を追加後、100問版に拡張予定です。

\mondaipage
\begin{honshikenmoshi}{本試験形式模試 第1回}

制限時間の目安: \textbf{25分}（20問 × 約75秒）

\vspace{0.5em}
\noindent\textbf{【ストラテジ系】}
\vspace{0.3em}

% --- Q1 ストラテジ ---
\begin{mondai}{CSR（Corporate Social Responsibility）の説明として最も適切なものはどれか。}{basic}
  \choice{企業が利益を最大化するための戦略}
  \choice{企業が社会的責任を果たすための活動}
  \choice{企業の内部統制を整備する仕組み}
  \choice{企業が従業員に対して行う福利厚生}
\end{mondai}

% --- Q2 ストラテジ ---
\begin{mondai}{RPA（Robotic Process Automation）の導入効果として最も適切なものはどれか。}{standard}
  \choice{工場の生産ラインをロボットで自動化する}
  \choice{定型的なPC操作をソフトウェアで自動化して業務効率を向上させる}
  \choice{AI が人間に代わって経営判断を行う}
  \choice{すべての業務プロセスを完全に無人化する}
\end{mondai}

% --- Q3 ストラテジ ---
\begin{mondai}{コンプライアンスの説明として最も適切なものはどれか。}{basic}
  \choice{業務プロセスの改善活動}
  \choice{法令や社会規範を遵守すること}
  \choice{情報システムの外部委託}
  \choice{新規事業の企画立案}
\end{mondai}

% --- Q4 ストラテジ ---
\begin{mondai}{IoT（Internet of Things）の事例として最も適切なものはどれか。}{standard}
  \choice{社内のファイルサーバをクラウドストレージに移行する}
  \choice{工場の機械にセンサを取り付け、稼働状況をリアルタイムで監視する}
  \choice{社員全員にスマートフォンを支給する}
  \choice{社内の業務システムをWebアプリケーションに置き換える}
\end{mondai}

% --- Q5 ストラテジ ---
\begin{mondai}{ABC分析の説明として最も適切なものはどれか。}{standard}
  \choice{売上高や在庫金額の大きい順に並べ、A・B・Cの3グループに分類して管理する手法}
  \choice{製品を花形・金のなる木・問題児・負け犬に分類する手法}
  \choice{顧客を年齢・性別・地域で3つに分類する手法}
  \choice{品質不良の原因を特性要因図で分類する手法}
\end{mondai}

% --- Q6 ストラテジ ---
\begin{mondai}{不正競争防止法で保護される「営業秘密」の3要件に含まれないものはどれか。}{trick}
  \choice{秘密管理性}
  \choice{有用性}
  \choice{非公知性}
  \choice{独創性}
\end{mondai}

% --- Q7 ストラテジ ---
\begin{mondai}{BPR（Business Process Reengineering）の説明として最も適切なものはどれか。}{standard}
  \choice{既存の業務プロセスを小さく改善し続けること}
  \choice{業務プロセスを根本的に再設計し、劇的な改善を図ること}
  \choice{業務プロセスをそのまま情報システムに置き換えること}
  \choice{外部企業に業務プロセスを委託すること}
\end{mondai}

\vspace{0.5em}
\noindent\textbf{【マネジメント系】}
\vspace{0.3em}

% --- Q8 マネジメント ---
\begin{mondai}{プロジェクトスコープの説明として最も適切なものはどれか。}{basic}
  \choice{プロジェクトの予算と費用}
  \choice{プロジェクトで実現すべき成果物や作業の範囲}
  \choice{プロジェクトの開始日と終了日}
  \choice{プロジェクトに参加するメンバーの人数}
\end{mondai}

% --- Q9 マネジメント ---
\begin{mondai}{ガントチャートの説明として最も適切なものはどれか。}{basic}
  \choice{作業の依存関係をネットワーク図で示したもの}
  \choice{作業のスケジュールを横棒（バー）で時系列に示したもの}
  \choice{作業の費用を円グラフで示したもの}
  \choice{品質の推移を折れ線グラフで示したもの}
\end{mondai}

% --- Q10 マネジメント ---
\begin{mondai}{ウォーターフォール開発モデルの特徴として最も適切なものはどれか。}{standard}
  \choice{要件定義・設計・実装・テスト・運用の各工程を順に進める}
  \choice{短い反復サイクルで機能を段階的にリリースする}
  \choice{プロトタイプを繰り返し作成して仕様を確定する}
  \choice{開発作業を外部に委託し、完成品を検収する}
\end{mondai}

% --- Q11 マネジメント ---
\begin{mondai}{ITサービスマネジメントにおける「問題管理」と「インシデント管理」の違いとして最も適切なものはどれか。}{trick}
  \choice{インシデント管理は根本原因を除去し、問題管理はサービスを復旧する}
  \choice{問題管理は根本原因を究明して再発防止を図り、インシデント管理はサービスの迅速な復旧を図る}
  \choice{両者は同じ活動の別名である}
  \choice{問題管理は利用者が行い、インシデント管理はサービス提供者が行う}
\end{mondai}

% --- Q12 マネジメント ---
\begin{mondai}{共通フレームにおいて、システム開発の最初の工程はどれか。}{standard}
  \choice{システム設計}
  \choice{要件定義}
  \choice{プログラミング}
  \choice{システムテスト}
\end{mondai}

\vspace{0.5em}
\noindent\textbf{【テクノロジ系】}
\vspace{0.3em}

% --- Q13 テクノロジ ---
\begin{mondai}{OSS（オープンソースソフトウェア）の特徴として最も適切なものはどれか。}{basic}
  \choice{無料で利用できるが、ソースコードは公開されていない}
  \choice{ソースコードが公開され、自由に利用・改変・再配布できる}
  \choice{特定の企業だけが利用を許可されたソフトウェアである}
  \choice{ソースコードは公開されているが、改変は禁止されている}
\end{mondai}

% --- Q14 テクノロジ ---
\begin{mondai}{SQLのSELECT文で、条件を指定して行を抽出する句はどれか。}{basic}
  \choice{FROM句}
  \choice{ORDER BY句}
  \choice{WHERE句}
  \choice{GROUP BY句}
\end{mondai}

% --- Q15 テクノロジ ---
\begin{mondai}{VPN（Virtual Private Network）の説明として最も適切なものはどれか。}{standard}
  \choice{社内LANの通信速度を高速化する技術}
  \choice{公衆回線上に仮想的な専用ネットワークを構築する技術}
  \choice{Webサイトのアクセスを制限する技術}
  \choice{無線LAN の電波到達範囲を広げる技術}
\end{mondai}

% --- Q16 テクノロジ ---
\begin{mondai}{ランサムウェアの説明として最も適切なものはどれか。}{standard}
  \choice{キーボードの入力内容を記録して外部に送信するソフトウェア}
  \choice{コンピュータのファイルを暗号化し、復号のための身代金を要求するマルウェア}
  \choice{正規のソフトウェアを装い、裏で不正な動作をするソフトウェア}
  \choice{ネットワーク上の他のコンピュータに自動的に感染を広げるソフトウェア}
\end{mondai}

% --- Q17 テクノロジ ---
\begin{mondai}{2進数 1011 を10進数に変換した値はどれか。}{basic}
  \choice{9}
  \choice{10}
  \choice{11}
  \choice{13}
\end{mondai}

% --- Q18 テクノロジ ---
\begin{mondai}{DBMS（DataBase Management System）が提供する機能として\textbf{適切でないもの}はどれか。}{trick}
  \choice{データの整合性を維持する排他制御}
  \choice{障害時にデータを復旧するリカバリ機能}
  \choice{ネットワーク上の通信経路を選択するルーティング機能}
  \choice{利用者のアクセスを制御するセキュリティ機能}
\end{mondai}

% --- Q19 テクノロジ ---
\begin{mondai}{クラウドサービスにおいて、サービス提供者がアプリケーションまで管理し、利用者はアプリケーションの機能だけを利用する形態はどれか。}{standard}
  \choice{IaaS}
  \choice{PaaS}
  \choice{SaaS}
  \choice{DaaS}
\end{mondai}

% --- Q20 テクノロジ ---
\begin{mondai}{多要素認証の組合せとして適切なものはどれか。}{trick}
  \choice{パスワードとPINコード}
  \choice{ICカードと指紋認証}
  \choice{2種類のパスワードの入力}
  \choice{秘密の質問と合い言葉}
\end{mondai}

\end{honshikenmoshi}

% -----------------------------------------------
%  解答・解説
% -----------------------------------------------
\kaitoupage
\subsection*{本試験形式模試 第1回 解答・解説}

\noindent\textbf{【ストラテジ系】}

% Q1
\answer{B}
\begin{kaisetsu}
\textbf{問1:} CSR は Corporate Social Responsibility（企業の社会的責任）。環境保全、地域貢献、法令遵守など、利益追求だけでなく社会全体に対する責任を果たす活動のこと。A は利益最大化戦略の一般論、C は内部統制、D は福利厚生であり、いずれもCSRの定義とは異なる。
\end{kaisetsu}

% Q2
\answer{B}
\begin{kaisetsu}
\textbf{問2:} RPA は「定型的なPC作業をソフトウェアのロボットで自動化する」技術。工場のロボット（A）とは異なり、PC上の入力・転記・集計などの事務作業が対象。C のようにAIが経営判断をするわけではなく、D のように全業務を無人化するわけでもない。
\end{kaisetsu}

% Q3
\answer{B}
\begin{kaisetsu}
\textbf{問3:} コンプライアンスは「法令遵守」。法律だけでなく社内規則や社会的倫理も含む。企業不祥事の防止に直結する重要概念。
\end{kaisetsu}

% Q4
\answer{B}
\begin{kaisetsu}
\textbf{問4:} IoTは「モノのインターネット」。センサ付き機器がインターネットに接続され、データの収集・活用を行うのが特徴。A はクラウド移行、C はモバイル配備、D はWebアプリ化であり、いずれも「モノがインターネットにつながる」話ではない。
\end{kaisetsu}

% Q5
\answer{A}
\begin{kaisetsu}
\textbf{問5:} ABC分析はパレート図を用いて、売上高等を大きい順に並べてA（上位70\%）・B（次の20\%）・C（残り10\%）のように分類し、重点管理する手法。B はPPM、D は特性要因図（フィッシュボーン図）。
\end{kaisetsu}

% Q6
\answer{D}
\begin{kaisetsu}
\textbf{問6:} 営業秘密の3要件は「秘密管理性」「有用性」「非公知性」。「独創性」は含まれない（特許の要件に近い）。この3つがすべて揃わないと不正競争防止法で保護されない。
\end{kaisetsu}

% Q7
\answer{B}
\begin{kaisetsu}
\textbf{問7:} BPRは業務を\textbf{抜本的に再設計}すること。A の漸進的改善はBPI（Business Process Improvement）やカイゼン。BPRは「ゼロベース」で業務を見直す点が特徴。
\end{kaisetsu}

\vspace{0.5em}
\noindent\textbf{【マネジメント系】}

% Q8
\answer{B}
\begin{kaisetsu}
\textbf{問8:} プロジェクトスコープは「何を作り、何をやるか」の範囲。スコープが曖昧だと手戻りが発生する。A はコスト、C はスケジュール、D は資源管理の話。
\end{kaisetsu}

% Q9
\answer{B}
\begin{kaisetsu}
\textbf{問9:} ガントチャートは横軸に時間、縦軸に作業を取り、各作業の期間をバー（横棒）で表す。A のネットワーク図はPERT図（アローダイアグラム）の説明。
\end{kaisetsu}

% Q10
\answer{A}
\begin{kaisetsu}
\textbf{問10:} ウォーターフォールモデルは各工程を順に進め、原則として前の工程に戻らない。B はアジャイル開発、C はプロトタイピングモデルの特徴。
\end{kaisetsu}

% Q11
\answer{B}
\begin{kaisetsu}
\textbf{問11:} インシデント管理＝「まず復旧（応急処置）」、問題管理＝「根本原因を突き止めて再発防止（根治）」。A は逆になっている。この2つの違いは頻出。
\end{kaisetsu}

% Q12
\answer{B}
\begin{kaisetsu}
\textbf{問12:} 共通フレームでは、まず「要件定義」で利用者のニーズを明確にする。その後に設計→実装→テストと進む。要件定義を飛ばして設計を始めると手戻りが起きる。
\end{kaisetsu}

\vspace{0.5em}
\noindent\textbf{【テクノロジ系】}

% Q13
\answer{B}
\begin{kaisetsu}
\textbf{問13:} OSSはソースコードが公開され、ライセンスに従って自由に利用・改変・再配布が可能。Linux、Apache、PostgreSQL などが代表例。「無料＝OSS」ではなく「ソースコード公開」が本質。
\end{kaisetsu}

% Q14
\answer{C}
\begin{kaisetsu}
\textbf{問14:} WHERE句は抽出条件を指定する。FROM句は対象の表、ORDER BY句は並び替え、GROUP BY句はグループ化。SELECT〜FROM〜WHERE〜の基本構文を覚えておくこと。
\end{kaisetsu}

% Q15
\answer{B}
\begin{kaisetsu}
\textbf{問15:} VPNはインターネット等の公衆回線上に暗号化技術を使って仮想的な専用線を構築する。テレワークやリモートアクセスに広く使われている。
\end{kaisetsu}

% Q16
\answer{B}
\begin{kaisetsu}
\textbf{問16:} ランサムウェアは「身代金（ransom）」を要求するマルウェア。ファイルを暗号化して使えなくし、復号と引き換えに金銭を要求する。A はキーロガー、C はトロイの木馬、D はワーム。
\end{kaisetsu}

% Q17
\answer{C}
\begin{kaisetsu}
\textbf{問17:} 2進数 1011 = $1 \times 2^3 + 0 \times 2^2 + 1 \times 2^1 + 1 \times 2^0 = 8 + 0 + 2 + 1 = 11$。各桁を右から $2^0, 2^1, 2^2, 2^3$ の重みで計算する。
\end{kaisetsu}

% Q18
\answer{C}
\begin{kaisetsu}
\textbf{問18:} ルーティングはルータやL3スイッチの機能であり、DBMSの機能ではない。DBMSは排他制御（A）、リカバリ（B）、アクセス制御（D）のほか、トランザクション管理などを提供する。
\end{kaisetsu}

% Q19
\answer{C}
\begin{kaisetsu}
\textbf{問19:} SaaS はアプリケーションを含めてすべてをサービス提供者が管理する。利用者はブラウザ等でアプリケーション機能だけを利用する。IaaS（A）はインフラのみ、PaaS（B）は開発基盤まで。
\end{kaisetsu}

% Q20
\answer{B}
\begin{kaisetsu}
\textbf{問20:} 多要素認証は「知識」「所持」「生体」の\textbf{異なる種類}を2つ以上組み合わせる。ICカード（所持）＋指紋（生体）は異なる2要素。A のパスワードとPINはどちらも「知識」、C も「知識」×2、D も「知識」×2 であり、多要素ではなく多段階認証にすぎない。
\end{kaisetsu}

% ====================================================
%  6.3 公式問題セクション（受け皿）
% ====================================================
\section{公式問題セクション（受け皿）}

このセクションは、IPA公式の過去問題を追加するための\textbf{受け皿}です。

\begin{keypoint}
\textbf{公式問題追加の方針:}
\begin{itemize}
  \item IPAが公開している過去問から、頻出テーマ・良問を厳選して収録予定
  \item 各問題には出典（試験回・問題番号）を明記
  \item 解説はオリジナル問題と同じフォーマットで統一
  \item 本試験形式模試（6.2節）を100問版に拡張する際にも活用
\end{itemize}
\end{keypoint}

\vspace{1em}
\noindent 以下に、公式問題を追加する際のフォーマット例を示します。

\vspace{0.5em}

\noindent\textbf{【フォーマット例】}

\begin{tcolorbox}[
  enhanced,
  colback=LightGray,
  colframe=SubBlue,
  boxrule=0.4pt,
  arc=3pt,
  left=8pt, right=8pt, top=6pt, bottom=6pt
]
{\small
\begin{verbatim}
\begin{mondai}{問題文をここに記載する。}{standard}
  \choice{選択肢ア}
  \choice{選択肢イ}
  \choice{選択肢ウ}
  \choice{選択肢エ}
\end{mondai}
\officialtag{R6}{42}

\answer{B}

\begin{kaisetsu}
解説をここに記載する。
\end{kaisetsu}
\end{verbatim}
}
\end{tcolorbox}

\vspace{0.5em}
\noindent 上記の |\officialtag{R6}{42}| は「令和6年 問42」を意味します。出典を明記することで、公式問題とオリジナル問題を明確に区別します。

\vspace{1em}
\noindent （公式過去問は今後のアップデートで追加予定です。）

% ====================================================
%  6.4 最終チェックリスト
% ====================================================
\section{最終チェックリスト}

試験直前に確認しましょう。以下の項目について「自分の言葉で説明できるか」をチェックしてください。説明できない項目があれば、該当する章に戻って復習しましょう。

\vspace{0.5em}

\subsection*{ストラテジ系の重要テーマ}

\begin{itemize}
  \item[$\square$] SWOT分析の4要素（内部/外部 × プラス/マイナス）を説明できる
  \item[$\square$] 3C分析の3つのCが何か言える
  \item[$\square$] PPMの4象限と2つの軸を説明できる
  \item[$\square$] バリューチェーンの主活動と支援活動の違いが分かる
  \item[$\square$] マーケティングの4Pを言える（Placeは「流通」）
  \item[$\square$] PL・BS・CFの違いを説明できる
  \item[$\square$] 請負・準委任・派遣の違い（特に指揮命令権）を区別できる
  \item[$\square$] 著作権・特許権・商標権の違い（登録要否・保護期間）を覚えている
  \item[$\square$] DX・IoT・AI・RPAの違いを説明できる
  \item[$\square$] SaaS・PaaS・IaaSの違いを覚えている
\end{itemize}

\subsection*{マネジメント系の重要テーマ}

\begin{itemize}
  \item[$\square$] WBSの目的と作り方が分かる
  \item[$\square$] ウォーターフォールとアジャイルの違いを説明できる
  \item[$\square$] ガントチャートとアローダイアグラムの違いが分かる
  \item[$\square$] PMBOK の知識エリア（スコープ、コスト、スケジュール等）を理解している
  \item[$\square$] ITILのインシデント管理と問題管理の違いを区別できる
  \item[$\square$] SLAが何か説明できる
  \item[$\square$] システム監査の独立性を理解している
  \item[$\square$] 共通フレームの工程の流れを説明できる
\end{itemize}

\subsection*{テクノロジ系の重要テーマ}

\begin{itemize}
  \item[$\square$] 情報セキュリティの3要素（CIA）を言える
  \item[$\square$] 公開鍵暗号と共通鍵暗号の違いを説明できる
  \item[$\square$] 多要素認証の「3つの要素」（知識・所持・生体）が分かる
  \item[$\square$] マルウェアの種類（ランサムウェア、トロイの木馬、ワーム等）を区別できる
  \item[$\square$] TCP/IPの4階層モデルと代表的なプロトコルを覚えている
  \item[$\square$] 関係データベースの主キー・外部キーの違いが分かる
  \item[$\square$] SQLの基本（SELECT, FROM, WHERE）を理解している
  \item[$\square$] 2進数と10進数の変換ができる
  \item[$\square$] VPN・ファイアウォール・DMZの役割を説明できる
  \item[$\square$] OSSの特徴を説明できる
\end{itemize}

\vspace{1em}

\begin{column}{試験当日のコツ}
\textbf{1. 時間配分を意識する}\\
100問で120分。1問あたり約72秒が目安です。分からない問題に時間をかけすぎず、マークして先に進みましょう。CBT方式なので「あとで見直す」フラグが使えます。

\vspace{0.5em}
\textbf{2. 消去法を活用する}\\
4択のうち「明らかに違う選択肢」を先に消せば、正解率が大きく上がります。2つに絞れたら50\%。「分からない」は「消去法で2択まで絞る」に変換しましょう。

\vspace{0.5em}
\textbf{3. 問題文の「最も適切なもの」「含まれないもの」に注意}\\
否定形の設問（「適切でないものはどれか」）を肯定形と読み間違えるミスは非常に多いです。設問のキーワードに下線を引く習慣をつけましょう。

\vspace{0.5em}
\textbf{4. 体調管理を忘れずに}\\
120分の集中は体力を消耗します。前日は十分な睡眠を取り、試験前の軽い食事と水分補給を忘れずに。会場には早めに到着して気持ちを落ち着けましょう。

\vspace{0.5em}
\textbf{5. 見直し時間を確保する}\\
最後の10〜15分を見直しに使えるよう、ペース配分を意識しましょう。マークミスや読み間違いの発見だけでも、数点の上乗せになります。
\end{column}

% ====================================================
%  章末QR
% ====================================================

\qrchapter{6}
